\chapter{A path attribution is not an entitlement}\label{ch:attribution}
Once the group value is known, one may want a number beside each child
action. A decomposition can be mathematically exact without being a uniquely
identified causal contribution. The difference must be kept visible before
any such number is connected to a persistent account.

\section{The declared common-path decomposition}
Choose the straight group path $\gamma(s)=z+s\delta_G$. Define
\begin{equation}\label{eq:attribution}
 R_a=-\int_0^1\mu(z+s\delta_G)^T\delta_a\,ds
     =-\mu(z)^T\delta_a-\tfrac12\delta_a^TH\delta_G.
\end{equation}
The last equality uses the affine Gaussian gradient. Sum over children:
\begin{align}
 \sum_aR_a
 &=-\int_0^1\mu(z+s\delta_G)^T\sum_a\delta_a\,ds\\
 &=-\int_0^1\mu(z+s\delta_G)^T\delta_G\,ds=E_G.
\end{align}
This proves closure for the cost-free domain. If $C_G\ne0$, the path
terms still sum to the field component; allocating or accounting for cost
requires its own declaration. It must not disappear in the notation.

\section{What the formula assigns in a small group}
Let the children send one and three units from Ridge to Vale at $(18,2)$.
The group increment is $(-4,4)$ and its value is twelve. Initial tangent
benefits are four and twelve. Curvature corrections are one and three, so
$R_A=3$, $R_B=9$. The total closes exactly.

Their standalone values are $15/4$ and $39/4$, summing to $27/2$.
Their sequential live values depend on order. Sending one first gives
$15/4$ followed by $33/4$; sending three first gives $39/4$ followed by
$9/4$. Both sequential sums equal twelve. Neither ordered split is the
same as $(3,9)$, because it describes another allocation convention.

The sum identity cannot tell us which of these child descriptions is fair,
deserved or causally identified. It tells us which declared arithmetic is
being used and whether that arithmetic closes on the group.

\section{Physical path versus interpolation}
Under a declared constant-rate physical model, each child's normalized
rate contributes $\delta_a$ along the common path. The integral then
decomposes that model path. Under endpoint interpolation it is an accounting
construction. Identical arithmetic does not erase the distinction.

Historical sources use the name radial Aumann--Shapley for related path
allocation \cite{foundation}. Using or avoiding that name does not resolve
the interpretation. The scientifically useful record specifies the actual
path and assumptions, labels the result as a mathematical attribution, and
states what it does not establish: ownership, fairness, causal entitlement,
or permission to spend.

\section{Splitting and cancellation}
Suppose one child is replaced by two proportional subchildren with increments
$\lambda\delta_a$ and $(1-\lambda)\delta_a$, while the group path is
unchanged. The formula is linear in the child increment, so their attributed
values add to the original $R_a$. This is a limited mathematical splitting
property. A split that changes grouping, fees, ownership or feasibility is
not covered by it.

If $\delta_G=0$, the formula becomes $R_a=-\mu(z)^T\delta_a$.
These child values can be nonzero and cancel. In the opposing two-unit
example they are $+8$ and $-8$, while the group has zero value. This
illustration makes the interpretive risk vivid: a cancelling endpoint does
not establish a gift of eight units to one participant and a legitimate
charge of eight to another.

\section{From children to owners is another map}
If a prospective policy assigns every child a unique owner, then finite-sum
rearrangement gives
\begin{equation}
 \Delta B_i=\sum_{a:\,\operatorname{owner}(a)=i}R_a,
 \qquad \sum_i\Delta B_i=\sum_aR_a.
\end{equation}
The second equality is exact once the map exists. It does not select the
map. A requester, provider, beneficiary and physical source can be different
entities. Substituting one for another because a software field is convenient
would change the policy while leaving the arithmetic looking correct.

\section*{Chapter recap}
Common-path attribution provides an explicit, auditable decomposition. Its
proof closes a sum. Ownership and spendability require separate prospective
authority even for a group with only one child.
